% Part: methods
% Chapter: proofs
% Section: proof-by-contradiction

\documentclass[../../../include/open-logic-section]{subfiles}

\begin{document}

\olfileid{mth}{prf}{con}

\olsection{反证法}

首先，反证法是一种用来证明否定性主张的推理模式。假设你想证明某个主张~$p$ 是\emph{假的}，即你想证明~$\lnot p$。最可行的策略是： 假设 $p$~为真，然后 证明这一假设会导致某个已知为假的结论。这个“已知为假的结论”可能是与~$p$ 本身相冲突（即相矛盾）的结果，也可能是与你所考虑的整体主张的其他假设相矛盾的结果。例如，证明“如果 $q$ 则 $\lnot p$”就涉及假设 $q$~为真并从中证明~$\lnot p$。如果你用反证法来证明 $\lnot p$，那就意味着除了~$q$ 之外还要假设~$p$。如果你能从 $p$ 推出 $\lnot q$，那么你就表明了假设~$p$ 会导致与你另一个假设~$q$ 相矛盾的结论，因为 $q$~和 $\lnot q$ 不能同时为真。当然，在证明矛盾的过程中，你还需要用到其他推理模式，以及展开定义。让我们来看一个例子。

\begin{prop}
  如果 $A \subseteq B$ 且 $B = \emptyset$，那么 $A$ 没有元素。
\end{prop}

\begin{proof}
  假设 $A \subseteq B$ 且 $B = \emptyset$。我们想要证明 $A$ 没有元素。
  \begin{quote}
    由于这是一个条件句，我们假设前件并想要证明后件。后件是：$A$ 没有元素。我们可以说得更明确一些：不存在~$x \in A$。
  \end{quote}
  $A$ 没有元素，当且仅当不存在~$x$ 使得 $x \in A$。
  \begin{quote}
    于是我们确定了要证明的实际上是一个否定主张~$\lnot p$，即：不存在 $x \in A$。为了使用反证法，我们必须假设相应的肯定主张~$p$，即存在 $x \in A$，并从中推出矛盾。我们通常写“用反证法，假设”或干脆写“假设不然”，然后陈述假设~$p$，以此标明我们在使用反证法。
  \end{quote}
  假设不然：存在 $x \in A$。
  \begin{quote}
    这是我们现在将用来推出矛盾的新假设。我们还有两个假设：$A \subseteq B$ 和 $B = \emptyset$。第一个假设给出 $x \in B$：
  \end{quote}
  因为 $A \subseteq B$，所以 $x \in B$。
  \begin{quote}
    但由于 $B = \emptyset$，$B$ 的每个元素（例如 $x$）也必须是~$\emptyset$ 的元素。
  \end{quote}
  因为 $B = \emptyset$，所以 $x \in \emptyset$。这是一个矛盾，因为根据定义 $\emptyset$ 没有元素。
  \begin{quote}
    这已经完成了证明：我们从所作的假设得出了所需的结论（一个矛盾），这意味着这些假设不可能全部成立。由于前两个假设（$A \subseteq B$ 和 $B = \emptyset$）不被质疑，那么一定是最后引入的那个假设（存在 $x \in A$）必定为假。但如果我们想做得彻底，可以把它明确说出来。
  \end{quote}
  因此，存在 $x \in A$ 这一假设必定为假，于是，根据反证法，$A$ 没有元素。
\end{proof}

每个肯定论断都平凡地等价于一个否定论断：$p$ 当且仅当 $\lnot\lnot p$。
因此，反证法也能用来“间接地”确立肯定论断，如下：要证明~$p$，可以把它视作否定论断 $\lnot\lnot p$。
如果我们能从 $\lnot p$ 推出矛盾，那么就通过反证法确立了 $\lnot\lnot p$，从而得到~$p$。

在之前的例子中，我们的目标是证明一个否定论断，即 $A$ 没有元素，因此我们为反证法所作的假设（即存在 $x \in A$）是一个肯定论断。
它为我们提供了可操作的对象，即那个假想的 $x \in A$，我们围绕它继续推理，直到得出 $x \in \emptyset$。

当间接证明一个肯定论断时，你为反证法所作的假设将是否定的。
但你常常可以很容易地将一个肯定论断重新表述为否定论断，反之亦然。
如果我们当初证明的是“$A = \emptyset$”而非那个否定性的结论“$A$ 没有元素”，
我们之前的证明本质上会是相同的。（根据 $=$ 的定义，“$A = \emptyset$”是一个普遍论断，因为它展开来就是“$A$ 的每个元素都是 $\emptyset$ 的元素，反之亦然”。)
但很容易看出它等价于否定论断“并非：存在 $x \in A$”。

所以，有时将 $\lnot p$ 作为假设来处理，比直接证明~$p$ 更容易。
即使直接证明同样简单甚至更简单（如下一例），有些人也倾向于使用间接证明。
如果双重否定让你困惑，不妨把对某个论断的反证法看作是从其\emph{相反}论断出发推出矛盾的证明。
因此，对 $\lnot p$ 的反证法就是从假设~$p$ 出发推出矛盾；而对~$p$ 的反证法则是从~$\lnot p$ 出发推出矛盾。

\begin{prop}
$A \subseteq A \cup B$。
\end{prop}

\begin{proof}
  我们要证明 $A \subseteq A \cup B$。
  \begin{quote}
    表面上看，这是一个肯定论断：每个 $x \in A$ 也在 $A \cup B$ 中。
    它的否定是：某个 $x \in A$ 不在 $A \cup B$ 中。
    因此，我们可以通过假设这个被否定的论断，并证明它导致矛盾，来间接地证明原论断。
  \end{quote}
  假设不然，即 $A \nsubseteq A \cup B$。
  \begin{quote}
    我们有 $A \subseteq A \cup B$ 的定义：每个 $x \in A$ 也属于 $A \cup B$。
    要理解 $A \nsubseteq A \cup B$ 意味着什么，我们需要对这个展开的定义做一些基本的逻辑处理：
    “每个 $x \in A$ 也属于 $A \cup B$”为假，当且仅当存在\emph{某个}~$x \in A$ 不属于 $A \cup B$。
    （这里你需要非常小心：许多学生尝试的反证法之所以失败，正是因为他们错误地分析了像“所有 $A$ 都是 $B$”这类论断的否定。）
    换言之，$A \nsubseteq A \cup B$ 当且仅当存在 $x$ 使得 $x \in A$ 且 $x \notin A \cup B$。从此往后，就容易了。
  \end{quote}
  因此，存在 $x \in A$ 使得 $x \notin A \cup B$。
  根据 $\cup$ 的定义，$x \in A \cup B$ 当且仅当 $x \in A$ 或 $x \in B$。
  由于 $x \in A$，我们有 $x \in A \cup B$。这与 $x \notin A \cup B$ 的假设矛盾。 
\end{proof}

\begin{prob}
\emph{间接地}证明 $A \cap B \subseteq A$。
\end{prob}

\begin{prop}
若 $A \subseteq B$ 且 $B \subseteq C$，则 $A \subseteq C$。
\end{prop}

\begin{proof}
  假设 $A \subseteq B$ 且 $B \subseteq C$。我们想证明 $A \subseteq C$。
  \begin{quote}
    让我们间接地进行：假设我们想要确立的结论的否定。
  \end{quote}
  假设不然，即 $A \nsubseteq C$。
  \begin{quote}
    如前所述，我们推理得出 $A \nsubseteq C$ 当且仅当并非每个 $x \in A$ 也属于 $C$，
    即存在某个 $x \in A$ 不属于 $C$。
    别担心，多练习后你就不必费神去分析这类否定了。
  \end{quote}
  换句话说，存在~$x$ 使得 $x \in A$ 且 $x \notin C$。
  \begin{quote}
    现在我们可以利用这点来得出矛盾。当然，我们需要使用其余两个假设来完成。
  \end{quote}
  由于 $A \subseteq B$，有 $x \in B$。由于 $B \subseteq C$，有 $x \in C$。但这与 $x \notin C$ 矛盾。
\end{proof}

\begin{prop}
若 $A \cup B = A \cap B$，则 $A = B$。
\end{prop}

\begin{proof}
  假设 $A \cup B = A \cap B$。我们要证 $A = B$。
  \begin{quote}
    现在开始常规步骤：
  \end{quote}
  用反证法，假设 $A \neq B$。
  \begin{quote}
    我们在反证法中所用的假设是 $A \neq B$。由于 $A = B$ 当且仅当 $A \subseteq B$ 且 $B \subseteq A$，可得 $A \neq B$ 当且仅当 $A \nsubseteq B$ \emph{或者} $B \nsubseteq A$。（注意在处理否定时小心谨慎有多么重要！{}）为了从这个析取式推出矛盾，我们将使用分情况证明，并证明在每种情况下都能推出矛盾。
  \end{quote}
  $A \neq B$ 当且仅当 $A \nsubseteq B$ 或者 $B \nsubseteq A$。我们分情况讨论。
  \begin{quote}
    在第一种情况下，我们假设 $A \nsubseteq B$，即存在某个 $x$，$x \in A$ 但 $x \notin B$。$A \cap B$ 定义为 $A$ 和 $B$ 的公共元素，所以只要某个元素不在其中一个集合中，它就不会在交集中。$A \cup B$ 是 $A$ 与 $B$ 的并集，因此属于任一集合的元素也在并集中。这告诉我们 $x \in A \cup B$ 但 $x \notin A \cap B$，从而 $A \cap B \neq A \cup B$。
  \end{quote}
  
  情况 1：$A \nsubseteq B$。则存在某个 $x$，使得 $x \in A$ 但 $x \notin B$。由于 $x \notin B$，则 $x \notin A \cap B$。由于 $x \in A$，$x \in A \cup B$。所以，$A \cap B \neq A \cup B$，这与假设 $A \cap B = A \cup B$ 矛盾。

  情况 2：$B \nsubseteq A$。则存在某个 $y$，使得 $y \in B$ 但 $y \notin A$。如前所述，我们有 $y \in A \cup B$ 但 $y \notin A \cap B$，因此 $A \cap B \neq A \cup B$，同样与 $A \cap B = A \cup B$ 矛盾。
\end{proof}

\end{document}